%==============================================================================
%  Explainable Machine Learning 2025/2026 -- University of Warsaw
%  Final report
%
%  Self-contained article-class file styled to match the JMLR course template.
%  Compiles out-of-the-box on Overleaf (pdfLaTeX). If you have the official
%  jmlr2e.sty from the course, you can port the body text into it; the section
%  structure here already follows the template (Intro / Methodology / Experiments
%  / Conclusions / References / Appendix).
%
%  Figures are left as TODO comments -- drop your screenshots where indicated.
%==============================================================================
\documentclass[11pt]{article}

\usepackage[utf8]{inputenc}
\usepackage[T1]{fontenc}
\usepackage{lmodern}
\usepackage[a4paper,margin=1in]{geometry}
\usepackage{microtype}
\usepackage{amsmath,amssymb}
\usepackage{booktabs}
\usepackage{graphicx}
\usepackage{float}        % provides [H] to pin tables/figures exactly in place
\usepackage{subcaption}
\usepackage{xcolor}
\usepackage{enumitem}
\usepackage[round]{natbib}     % author-year citations, JMLR style
\usepackage{fancyhdr}
\usepackage{titlesec}
\usepackage[hidelinks]{hyperref}

% ----- JMLR-like running heads -------------------------------------------------
\newcommand{\ShortAuthors}{Vinichenko}
\newcommand{\ShortTitle}{Causal Calibration of Saliency Maps}

\pagestyle{fancy}
\fancyhf{}
\fancyhead[L]{\footnotesize\ShortAuthors}
\fancyhead[R]{\footnotesize\ShortTitle}
\fancyfoot[C]{\thepage}
\renewcommand{\headrulewidth}{0pt}

\fancypagestyle{firstpage}{%
  \fancyhf{}
  \fancyhead[L]{\footnotesize Explainable Machine Learning 2025/2026 course at the University of Warsaw}
  \fancyhead[R]{\footnotesize Final report}
  \fancyfoot[C]{\thepage}
  \renewcommand{\headrulewidth}{0pt}
}

% ----- Section style (bold, numbered like the template) ------------------------
\titleformat{\section}{\large\bfseries}{\thesection.}{0.5em}{}
\titleformat{\subsection}{\normalsize\bfseries}{\thesubsection}{0.5em}{}
\setlength{\parindent}{0pt}
\setlength{\parskip}{4pt}

\title{\bfseries Closing the Reliability Gap: Causal Calibration of\\
Saliency Maps via Perturbation-Aware Temperature Scaling\\ and Significance Bucketing}
\author{}
\date{}

\begin{document}

%==============================================================================
\thispagestyle{firstpage}

\begin{center}
  {\LARGE\bfseries Closing the Reliability Gap: Causal Calibration of\\[2pt]
  Saliency Maps via Perturbation-Aware Temperature Scaling\\[2pt]
  and Significance Bucketing}\\[1.2em]
  % TODO: replace author block with the final author list / emails
  {\large\bfseries Oleksii Vinichenko}\hfill {\textsc{oleksiivinichenko@gmail.com}}\\[2pt]
  % {\large\bfseries Author B}\hfill {\textsc{b.stud@uw.edu.pl}}\\[2pt]
  \textit{University of Warsaw, Poland}
\end{center}

\vspace{1em}

%==============================================================================
\begin{center}\textbf{Abstract}\end{center}
\begin{quote}
\small
The insertion/deletion tests of \citet{petsiuk2018rise} are the de-facto causal
yardstick for saliency maps: pixels are ranked by attributed importance and
progressively removed or revealed while the model's confidence is tracked. We
study a \emph{reliability gap} in this protocol: the perturbed images that the
test feeds to the network are heavily out-of-distribution (OOD), so the softmax
confidence that the curve is built from is itself untrustworthy, biasing the
area-under-curve (AUC) score and producing noisy, non-monotone curves. Following
the ReCalX framework~\citep{recalx2025}, we recalibrate each model with
\emph{perturbation-aware temperature scaling}: one temperature per perturbation
bin, fitted by negative-log-likelihood minimisation, applied without changing the
model's predictions. We further propose \emph{significance bucketing}, a
rank-based perturbation schedule that removes equal-count groups of pixels
instead of a fixed top-$n\%$, making the test scale-invariant and robust to
sparse saliency maps. Across four ImageNet architectures (ResNet50/101,
ViT-B/16, ViT-B/32) and eight attribution methods, calibration lowers
expected calibration error (ECE $0.28\!\to\!0.15$) and thresholded adaptive
calibration error (TACE $0.35\!\to\!0.16$) on ResNet50, significantly increases
deletion AUC, and makes the causal curves significantly more monotone (Spearman
$r$: $-0.82\!\to\!-0.94$, Wilcoxon $p<0.01$), with the largest gains on
perturbation-based attributions. Significance bucketing is in turn significantly
more monotone than fixed-step removal ($p<0.001$). Together the two ingredients
make insertion/deletion a more trustworthy XAI metric.
\end{quote}

%==============================================================================
\section{Introduction}

Attribution (saliency) methods explain an image classifier by assigning each
input pixel a score for how much it supports a target class. To decide whether
such an explanation is \emph{faithful}, the community relies on the causal
insertion/deletion test~\citep{petsiuk2018rise}: rank pixels by importance, then
delete (or insert) them most-important-first and watch the predicted
probability. A faithful map makes confidence collapse quickly under deletion
(low AUC) and recover quickly under insertion (high AUC).

This report investigates a subtle but consequential flaw in that recipe and how
to fix it. Our three highlights are:

\begin{itemize}[leftmargin=1.4em,itemsep=2pt,topsep=2pt]
  \item \textbf{The reliability gap is real and measurable.} Masking pixels
  pushes inputs off the data manifold; the classifier becomes miscalibrated
  (ECE up to $0.35$ at high perturbation), so the confidence curve underlying
  the AUC is not trustworthy.
  \item \textbf{Perturbation-aware temperature scaling closes much of the gap.}
  A single temperature per perturbation level, fitted post-hoc and leaving
  predictions unchanged, cuts ECE/TACE roughly in half and makes deletion curves
  significantly more monotone.
  \item \textbf{Rank-based ``significance bucketing'' beats fixed-step removal.}
  Removing equal-count importance buckets is scale-invariant, robust to sparse
  maps (IG, Saliency), and yields significantly more monotone curves than the
  classic top-$n\%$ schedule.
\end{itemize}

\textbf{Research questions.}
\textbf{(RQ1)} Does perturbation-aware recalibration restore reliable confidence
under insertion/deletion perturbations, and does this make the resulting causal
curves more stable (monotone) and their AUC less biased?
\textbf{(RQ2)} Does significance bucketing provide a more sensitive and monotone
causal evaluation than fixed-step top-$n\%$ removal, especially for skewed
saliency maps?

\textbf{Contributions.}
\textbf{(C1)} A reproducible pipeline integrating eight attribution methods,
perturbation-binned temperature scaling (ReCalX), and two families of
perturbation generators (fixed-step and rank-based bucketing) across four
ImageNet architectures.
\textbf{(C2)} An empirical study quantifying the reliability gap and showing that
recalibration and bucketing each independently and significantly improve the
trustworthiness of insertion/deletion evaluation.
This project is directly related to the XAI course: it concerns the evaluation
and improvement of post-hoc explanations for deep vision models.

%==============================================================================
\section{Methodology}

\subsection{Attribution methods: inputs, outputs, meaning}
An attribution method takes a trained model $f$, an input image
$x\in\mathbb{R}^{3\times224\times224}$ and a target class $c$, and returns a 2D
importance map $S\in\mathbb{R}^{224\times224}$ (channels summed) where $S_{ij}$
scores how much pixel $(i,j)$ supports $f_c(x)$. We use eight methods spanning
two families. \emph{Gradient/back-propagation} methods --- Saliency
\citep{simonyan2014deep}, Integrated Gradients (IG)
\citep{sundararajan2017axiomatic} and GradientSHAP \citep{lundberg2017unified}
--- read importance from gradients of the class score. \emph{Perturbation}
methods --- Occlusion \citep{zeiler2014visualizing}, LIME
\citep{ribeiro2016why}, Feature Ablation / Shapley Value Sampling, and RISE
\citep{petsiuk2018rise} --- probe $f$ with masked inputs and infer importance
from the output change.

\textbf{RISE} is our default explainer because it is black-box and aligns
naturally with the perturbation philosophy of the causal test. It samples $N$
low-resolution binary masks ($s\times s$ grid, on-probability $p_1$), bilinearly
upsamples and randomly shifts them to image size, multiplies each with $x$,
queries the model, and estimates importance as the probability-weighted average
of the masks: $S = \frac{1}{N\,p_1}\sum_{m} f_c(x\odot m)\,m$. We use
$N{=}5000,\,s{=}10,\,p_1{=}0.1$ for ResNet50 and $N{=}2000$ for the other
architectures (speed).

\subsection{The causal insertion/deletion pipeline}
Given $S$, pixels are sorted by importance. \emph{Deletion} replaces the
top-ranked pixels with a baseline (black) in steps and records $f_c$; a faithful
map yields a fast drop and low AUC. \emph{Insertion} starts from a Gaussian-blur
baseline and restores top pixels; a faithful map yields a fast rise and high
AUC. We compute the AUC of the confidence-vs-fraction curve
(\texttt{sklearn.metrics.auc}, trapezoidal), and two stability statistics: the
\emph{Spearman} rank correlation between perturbation level and confidence
(ideal $-1$ for deletion, $+1$ for insertion) and a \emph{monotonicity ratio}
(fraction of steps moving in the expected direction).

\subsection{The reliability gap and ReCalX calibration}
Heavily masked images lie far from the training distribution, so $f$'s softmax
confidence becomes unreliable --- typically over-confident --- and the causal
curve is biased and noisy. \citet{recalx2025} formalise this link between
\emph{uncertainty calibration} and perturbation-based explanation quality and
propose \emph{ReCalX}: recalibrate the model for the specific perturbations used,
without altering its predictions. We adopt a binned temperature-scaling
realisation. The perturbation level $\ell\in[0,1]$ is split into $B{=}10$ bins;
for each bin $b$ we collect raw logits over a calibration set and fit a scalar
temperature $T_b$ by minimising cross-entropy with L-BFGS-B (bounds
$[0.05,5]$). At evaluation, the wrapper \texttt{ReCalXModel} selects the bin for
the current step and returns $z/T_b$ (softmax applied afterwards). Because
$\arg\max_c z_c$ is invariant to dividing by $T_b>0$, calibrated and raw models
make identical predictions; only the confidence is corrected.

\subsection{Significance bucketing}
The classic schedule removes a fixed top-$n\%$ (e.g.\ $1\%$) per step. This is
fragile for sparse/skewed maps: when most pixels share the same (near-zero)
value, value-thresholding collapses them into one step. \emph{Significance
bucketing} instead sorts pixels by \emph{rank} and removes $N$ equal-count
groups, one per step. Being rank-based it is scale-invariant (so IG, Saliency
and RISE are treated identically) and its $N$ steps line up with the calibration
bins. We use $N{=}25$ buckets for evaluation and contrast against fixed-step
$1\%$ deletion.

\subsection{Datasets, models, and the ViT bonus}
We use an ImageNet validation subset (4{,}358 images available; 20 sampled per
run) with standard preprocessing (resize 256, center-crop 224, ImageNet
normalisation). Models are torchvision ImageNet-pretrained ResNet50, ResNet101
\citep{he2016deep}, ViT-B/16 and ViT-B/32 \citep{dosovitskiy2021image}; ResNet50
and ViT-B/16 agree on the top-1 class for 19/20 sampled images. As a theoretical
bonus we surveyed three attribution methods designed specifically for Vision
Transformers: \emph{(i)} \textbf{Attention Rollout}
\citep{abnar2020quantifying}, which multiplies (residual-augmented) attention
matrices across layers to trace token-to-patch influence; \emph{(ii)}
\textbf{Transformer Attribution / LRP relevance propagation}
\citep{chefer2021transformer}, a class-specific rule propagating relevance
through attention and skip connections; and \emph{(iii)} \textbf{Generic
Attention-model Explainability} \citep{chefer2021generic}, which combines
gradient-weighted attention across layers and generalises to multi-modal
transformers. These are natural drop-in explainers for the ViT branch of our
pipeline.

%==============================================================================
\section{Experiments}

All runs use $N{=}20$ images; calibration temperatures are fit jointly (batch)
unless stated. ECE uses 10 equal-width bins; TACE uses 15 adaptive bins with
threshold $0.01$. Significance is assessed with the paired Wilcoxon signed-rank
test ($\alpha{=}0.05$).

% TODO (Figure 1): example image + saliency maps (RISE/IG/LIME ...) side by side.
% Source: notebooks/03_calibration_test.ipynb, Section 1/5 (attribution overlays).

% TODO (Figure 2): Deletion & Insertion confidence curves with AUC, raw vs
% calibrated. Source: visualizations.visualize_tests (Section 2 of the notebook).

\subsection{Phase 2 -- Pre-calibration baseline}
Table~\ref{tab:baseline} reports per-method deletion AUC (lower = sharper
confidence drop = better localisation) on ResNet50. Gradient methods (IG,
GradientSHAP, Saliency) achieve the lowest deletion AUC, while LIME is highest.
These raw numbers, however, are confounded by the reliability gap (next
subsection): the curves they integrate are built on miscalibrated confidence.

\begin{table}[H]
\centering\small
\caption{Phase 2 baseline on ResNet50 (RISE-style deletion, $N{=}20$): mean
deletion AUC and raw ECE per attribution method. Lower deletion AUC is better.}
\label{tab:baseline}
\begin{tabular}{lcc}
\toprule
Method & Deletion AUC $\downarrow$ & Raw ECE \\
\midrule
Integrated Gradients & $0.068 \pm 0.065$ & $0.162$ \\
GradientSHAP         & $0.080 \pm 0.081$ & $0.186$ \\
Saliency             & $0.111 \pm 0.075$ & $0.211$ \\
RISE                 & $0.178 \pm 0.113$ & $0.281$ \\
Feature Ablation     & $0.205 \pm 0.135$ & $0.390$ \\
LIME                 & $0.380 \pm 0.136$ & $0.422$ \\
\bottomrule
\end{tabular}
\end{table}

\subsection{Phase 3 -- Calibrated causal analysis}
\textbf{The gap.} On ResNet50 with RISE + bucket deletion, raw ECE/TACE rise
sharply with perturbation, confirming OOD over-confidence. The fitted per-bin
temperatures increase monotonically with perturbation
(Table~\ref{tab:temps}): heavier masking requires more softening, exactly the
signature ReCalX predicts.

\begin{table}[H]
\centering\small
\caption{Batch-fitted per-bin temperatures (ResNet50, RISE). $T>1$ softens
over-confident logits; the monotone rise tracks the increasing severity of
perturbation.}
\label{tab:temps}
\begin{tabular}{lcccccccccc}
\toprule
Bin (level) & 0--.1 & .1--.2 & .2--.3 & .3--.4 & .4--.5 & .5--.6 & .6--.7 & .7--.8 & .8--.9 & .9--1 \\
\midrule
$T_b$ & 0.24 & 0.52 & 0.57 & 0.55 & 0.60 & 0.67 & 0.70 & 0.79 & 0.96 & 1.37 \\
\bottomrule
\end{tabular}
\end{table}

\textbf{Calibration reduces miscalibration.} On ResNet50 (RISE,
bucket-deletion, $N{=}20$) recalibration roughly halves both errors:
ECE $0.282\!\to\!0.148$ and TACE $0.350\!\to\!0.160$; 15/20 images improve
(mean $\Delta$ECE $+0.134$). The effect is largest where raw miscalibration is
largest (Table~\ref{tab:percal}): LIME ($\Delta$ECE $-0.328$), Feature Ablation
($-0.186$) and RISE ($-0.134$), and near-neutral for the already-sharp gradient
maps (IG $+0.003$, GradientSHAP $-0.010$).

\begin{table}[H]
\centering\small
\caption{Per-method calibration effect on ResNet50 (ECE before/after,
$N{=}20$). Negative $\Delta$ECE = improvement.}
\label{tab:percal}
\begin{tabular}{lccc}
\toprule
Method & ECE raw & ECE cal & $\Delta$ECE \\
\midrule
LIME                 & 0.422 & 0.094 & $-0.328$ \\
Feature Ablation     & 0.390 & 0.205 & $-0.186$ \\
RISE                 & 0.281 & 0.148 & $-0.134$ \\
Saliency             & 0.211 & 0.168 & $-0.043$ \\
GradientSHAP         & 0.186 & 0.176 & $-0.010$ \\
Integrated Gradients & 0.162 & 0.165 & $+0.003$ \\
\bottomrule
\end{tabular}
\end{table}

% TODO (Figure 3): ECE-vs-perturbation curves, base vs ReCalX.
% Source: visualizations.plot_ece_curves (Section 3).
% TODO (Figure 4): boxplots of per-image ECE raw vs calibrated (Section 3).

\textbf{Calibration improves AUC and monotonicity.} A Wilcoxon test
(Table~\ref{tab:wilcoxon}) shows deletion AUC increases significantly after
calibration for \emph{every} method ($p<0.001$). Curve monotonicity (Spearman
$r$, more negative = better for deletion) improves significantly for the noisy
perturbation-based maps --- RISE $-0.82\!\to\!-0.94$ ($p{=}0.001$), LIME
$-0.59\!\to\!-0.94$ ($p<0.001$), Feature Ablation $-0.48\!\to\!-0.81$
($p{=}0.006$) --- while the gradient maps were already near-monotone and change
little. This directly answers RQ1: recalibration makes the causal signal both
larger and steadier, most where it was least reliable.

\begin{table}[H]
\centering\small
\caption{Effect of calibration on ResNet50 (Wilcoxon signed-rank, $N{=}20$).
$\Delta$ values are medians; $*$ marks $p<0.05$.}
\label{tab:wilcoxon}
\begin{tabular}{lcccc}
\toprule
Method & $\Delta$AUC & $p$(AUC) & $\Delta$Spearman & $p$(Spear) \\
\midrule
RISE             & $+0.250$ & $0.000^{*}$ & $-0.064$ & $0.001^{*}$ \\
LIME             & $+0.340$ & $0.000^{*}$ & $-0.309$ & $0.000^{*}$ \\
Feature Ablation & $+0.309$ & $0.000^{*}$ & $-0.254$ & $0.006^{*}$ \\
Saliency         & $+0.178$ & $0.000^{*}$ & $-0.016$ & $0.812$ \\
Integrated Grad. & $+0.082$ & $0.000^{*}$ & $+0.003$ & $0.317$ \\
GradientSHAP     & $+0.099$ & $0.000^{*}$ & $+0.012$ & $0.059$ \\
\bottomrule
\end{tabular}
\end{table}

\textbf{Across architectures.} Calibration helps all four backbones
(Table~\ref{tab:arch}), with the largest AUC gain on ResNet50 and consistently
more-negative Spearman after calibration. ViTs are better calibrated to begin
with (raw ECE $0.20$ vs $0.28$ for ResNet50) and already produce the most
monotone raw curves, so their headroom is smaller --- but calibration still
improves every metric.

\begin{table}[H]
\centering\small
\caption{Architecture comparison (RISE, bucket-deletion, $N{=}20$): deletion AUC
and Spearman $r$ before/after calibration.}
\label{tab:arch}
\begin{tabular}{lcccc}
\toprule
Architecture & AUC raw & AUC cal & Spear raw & Spear cal \\
\midrule
ResNet50  & 0.178 & 0.401 & $-0.816$ & $-0.942$ \\
ResNet101 & 0.275 & 0.344 & $-0.833$ & $-0.903$ \\
ViT-B/16  & 0.207 & 0.244 & $-0.924$ & $-0.962$ \\
ViT-B/32  & 0.210 & 0.244 & $-0.901$ & $-0.936$ \\
\bottomrule
\end{tabular}
\end{table}

% TODO (Figure 5): architecture-comparison boxplots (Section 10).

\subsection{Phase 3 -- Significance bucketing vs fixed-step}
Comparing the two perturbation schedules on ResNet50 + RISE
(Table~\ref{tab:bucket}), significance bucketing produces more monotone curves
than fixed-step $1\%$ deletion, both raw and calibrated. The differences are
statistically significant: Wilcoxon $p{=}0.0002$ for Spearman $r$ and
$p{=}0.0009$ for the monotonicity ratio. Calibration ECE is comparable between
the two ($0.128$ for TopN vs $0.148$ for Bucket), so bucketing buys stability
without hurting calibration. This answers RQ2: rank-based bucketing is the more
sensitive and robust schedule, and (being scale-invariant) it is the safer
default for sparse attribution maps.

\begin{table}[H]
\centering\small
\caption{Significance bucketing (25 buckets) vs fixed-step TopN ($1\%$) on
ResNet50 + RISE. Spearman $r$ closer to $-1$ and higher monotonicity are better.}
\label{tab:bucket}
\begin{tabular}{lcccc}
\toprule
Schedule & Spear raw & Spear cal & Mono raw & Mono cal \\
\midrule
TopN ($1\%$)     & $-0.722$ & $-0.908$ & 0.557 & 0.583 \\
Bucket ($25$)    & $-0.816$ & $-0.942$ & 0.632 & 0.686 \\
\midrule
\multicolumn{5}{l}{\footnotesize Wilcoxon Bucket vs TopN: Spearman $p{=}0.0002^{*}$; Monotonicity $p{=}0.0009^{*}$.}\\
\bottomrule
\end{tabular}
\end{table}

% TODO (Figure 6): TopN vs Bucket monotonicity boxplot (Section 11).

%==============================================================================
\section{Conclusions}
The insertion/deletion ``reliability gap'' is, to a large extent, a calibration
artefact: under explainability-specific perturbations the classifier is pushed
OOD and its confidence --- the very quantity the AUC integrates --- becomes
untrustworthy. Perturbation-aware temperature scaling (ReCalX) addresses this
cheaply and post-hoc: it leaves predictions untouched, roughly halves ECE/TACE,
significantly increases deletion AUC, and makes the causal curves significantly
more monotone, with the biggest gains on the perturbation-based attributions and
on CNNs. Rank-based significance bucketing further stabilises the test and is
robust to sparse maps, significantly outperforming fixed-step removal in
monotonicity. \textbf{Take-away:} calibrating the model \emph{for the
perturbation it is about to see}, and removing pixels by rank rather than by a
fixed percentage, together turn insertion/deletion into a markedly more
trustworthy XAI metric.

\textbf{Limitations and future work.} Evaluations use $N{=}20$ images, so effect
sizes should be read as indicative; the masking baseline (black/blur) is itself
OOD, motivating the planned inpainting-based generator (currently a stub); and
single-image calibration overfits (temperatures hit the search bounds), so
batch fitting is preferred. Promising extensions include superpixel/patch-based
deletion aligned with ViT tokenisation and FID/OTDD distances to quantify how
far each perturbation pushes inputs off-manifold.

%==============================================================================
\bibliographystyle{plainnat}
\begin{thebibliography}{99}
\setlength{\itemsep}{1pt}

\bibitem[Abnar and Zuidema(2020)]{abnar2020quantifying}
Samira Abnar and Willem Zuidema.
\newblock Quantifying attention flow in transformers.
\newblock In \emph{ACL}, 2020.

\bibitem[Chefer et al.(2021a)]{chefer2021transformer}
Hila Chefer, Shir Gur, and Lior Wolf.
\newblock Transformer interpretability beyond attention visualization.
\newblock In \emph{CVPR}, 2021.

\bibitem[Chefer et al.(2021b)]{chefer2021generic}
Hila Chefer, Shir Gur, and Lior Wolf.
\newblock Generic attention-model explainability for interpreting bi-modal and
encoder-decoder transformers.
\newblock In \emph{ICCV}, 2021.

\bibitem[Dosovitskiy et al.(2021)]{dosovitskiy2021image}
Alexey Dosovitskiy et al.
\newblock An image is worth 16x16 words: Transformers for image recognition at
scale.
\newblock In \emph{ICLR}, 2021.

\bibitem[He et al.(2016)]{he2016deep}
Kaiming He, Xiangyu Zhang, Shaoqing Ren, and Jian Sun.
\newblock Deep residual learning for image recognition.
\newblock In \emph{CVPR}, 2016.

\bibitem[Lundberg and Lee(2017)]{lundberg2017unified}
Scott M. Lundberg and Su-In Lee.
\newblock A unified approach to interpreting model predictions.
\newblock In \emph{NeurIPS}, 2017.

\bibitem[Petsiuk et al.(2018)]{petsiuk2018rise}
Vitali Petsiuk, Abir Das, and Kate Saenko.
\newblock {RISE}: Randomized input sampling for explanation of black-box models.
\newblock In \emph{BMVC}, 2018.
\newblock arXiv:1806.07421.

% TODO: verify author names for the ReCalX paper (arXiv:2511.10439, 2025).
\bibitem[ReCalX(2025)]{recalx2025}
ReCalX authors.
\newblock Recalibrating models for reliable perturbation-based explanations
({ReCalX}).
\newblock arXiv:2511.10439, 2025.

\bibitem[Ribeiro et al.(2016)]{ribeiro2016why}
Marco Tulio Ribeiro, Sameer Singh, and Carlos Guestrin.
\newblock ``Why should I trust you?'': Explaining the predictions of any
classifier.
\newblock In \emph{KDD}, 2016.

\bibitem[Simonyan et al.(2014)]{simonyan2014deep}
Karen Simonyan, Andrea Vedaldi, and Andrew Zisserman.
\newblock Deep inside convolutional networks: Visualising image classification
models and saliency maps.
\newblock In \emph{ICLR Workshop}, 2014.

\bibitem[Sundararajan et al.(2017)]{sundararajan2017axiomatic}
Mukund Sundararajan, Ankur Taly, and Qiqi Yan.
\newblock Axiomatic attribution for deep networks.
\newblock In \emph{ICML}, 2017.

\bibitem[Zeiler and Fergus(2014)]{zeiler2014visualizing}
Matthew D. Zeiler and Rob Fergus.
\newblock Visualizing and understanding convolutional networks.
\newblock In \emph{ECCV}, 2014.

\end{thebibliography}

%==============================================================================
\appendix
\section{Appendix: Additional Results}

\textbf{A.1 Integrated-Gradients cross-architecture run.} Repeating the
architecture comparison with IG instead of RISE (Table~\ref{tab:arch_ig})
reproduces the same qualitative pattern: calibration raises deletion AUC and
makes Spearman $r$ more negative on all four backbones.

\begin{table}[H]
\centering\small
\caption{Architecture comparison with Integrated Gradients (bucket-deletion,
$N{=}20$).}
\label{tab:arch_ig}
\begin{tabular}{lcccc}
\toprule
Architecture & AUC raw & AUC cal & Spear raw & Spear cal \\
\midrule
ResNet50  & 0.081 & 0.163 & $-0.900$ & $-0.936$ \\
ResNet101 & 0.190 & 0.220 & $-0.879$ & $-0.909$ \\
ViT-B/16  & 0.202 & 0.236 & $-0.909$ & $-0.935$ \\
ViT-B/32  & 0.173 & 0.213 & $-0.954$ & $-0.969$ \\
\bottomrule
\end{tabular}
\end{table}

\textbf{A.2 TopN vs Bucket calibration error.} ECE/TACE before$\to$after
calibration: TopN($1\%$) ECE $0.278\!\to\!0.128$, TACE $0.303\!\to\!0.132$;
Bucket($25$) ECE $0.281\!\to\!0.148$, TACE $0.350\!\to\!0.160$. Both schedules
calibrate comparably well.

\textbf{A.3 Spearman/monotonicity per method (raw $\to$ cal), ResNet50 + RISE.}
Feature Ablation $-0.48\!\to\!-0.81$, LIME $-0.59\!\to\!-0.94$, RISE
$-0.82\!\to\!-0.94$, Saliency $-0.89\!\to\!-0.79$, GradientSHAP
$-0.93\!\to\!-0.79$, IG $-0.93\!\to\!-0.81$. Calibration helps the noisy
(perturbation-based) maps and is roughly neutral-to-slightly-negative for maps
that were already highly monotone.

\textbf{A.4 Normalisation sanity check.} Raw attribution ranges differ wildly
across methods (e.g.\ IG $[-1.01,1.14]$, RISE $[0.18,1.21]$, GradientSHAP
$[-1.34,1.52]$). Because the bucket generators use \texttt{argsort}, they are
rank-invariant and unaffected by scale --- one of the practical motivations for
significance bucketing.

% TODO (Appendix figures): add any extra boxplots/curves you want to keep here;
% the annex has no length limit.

\end{document}
